\begin{table}[htbp]
\centering
\caption{Phân định vai trò các thuật toán mật mã trong giao thức TLS 1.3}
\label{tab:tls_algorithms}
\small
\renewcommand{\arraystretch}{1.3}
\begin{tabularx}{\textwidth}{@{}>{\raggedright\arraybackslash}p{2.8cm} >{\raggedright\arraybackslash}p{3.2cm} >{\raggedright\arraybackslash}p{3.4cm} >{\raggedright\arraybackslash}X@{}}
\toprule
\textbf{Lớp chức năng} & \textbf{Thuật toán tiêu biểu} & \textbf{Khóa mật mã sử dụng} & \textbf{Mục tiêu an ninh đáp ứng} \\
\midrule
\textbf{Xác thực (Authentication)} & 
RSA-PSS, ECDSA (secp256r1), Ed25519. & 
Khóa bất đối xứng dài hạn ($PK, SK$) trong chứng thư X.509. & 
Xác thực nguồn gốc máy chủ, ngăn chặn tấn công mạo danh và người đứng giữa (MITM). \\
\midrule
\textbf{Thỏa thuận khóa (Key Exchange)} & 
ECDHE (X25519, secp256r1), DHE. & 
Cặp khóa bất đối xứng tạm thời (ephemeral), sinh mới mỗi phiên. & 
Thiết lập bí mật chung an toàn; bảo đảm Bí mật chuyển tiếp hoàn hảo (PFS). \\
\midrule
\textbf{Mã hóa dữ liệu (Record Layer)} & 
AEAD (AES-128-GCM, AES-256-GCM, ChaCha20-Poly1305). & 
Khóa đối xứng phiên (Session Key) suy dẫn qua HKDF. & 
Bảo đảm tính bí mật và toàn vẹn dữ liệu ứng dụng với tốc độ và hiệu năng tối ưu. \\
\bottomrule
\end{tabularx}
\end{table}
